\documentclass[11pt]{article}
\usepackage[utf8]{inputenc}
\usepackage{amsmath,amssymb}
\usepackage{authblk}
\usepackage[margin=1in]{geometry}
\usepackage[hidelinks]{hyperref}
\usepackage{xurl}
\usepackage{setspace}

\title{Deterministic RealQM Beta Decay:\\
A Phase-Coincidence Half-Life}
\author[1]{Claes Johnson}
\affil[1]{KTH Royal Institute of Technology, SE-100 44 Stockholm, Sweden \\
\texttt{claesjohnson@gmail.com}}
\date{\today}

\begin{document}
\maketitle

\begin{abstract}
Radioactive decay poses two separate questions: \emph{when} a given nucleus decays, and \emph{what carries}
the released energy and momentum. Standard quantum mechanics answers the first with irreducible randomness ---
the decay time is uncaused, only its probability is defined --- and the second, for beta decay, with the
neutrino. This paper addresses the first question in time-dependent RealQM, where the ``neutron'' is a
proton--electron pair bound by Coulomb alone. Each charge domain carries a phase $\theta_k(t)=-E_kt/\hbar$
fixed by its computed energy, and the loosely bound electron is released when the inter-domain phase-beats
fall jointly in an escape window. On the reading of de Broglie and Bohm the decay is then
\emph{deterministic}, the randomness epistemic (ignorance of the phases), and the exponential, memoryless
law emerges as the ensemble statistics of a coincidence rather than as a signature of fundamental
indeterminism. We are equally explicit that the specific coexisting-domain
mechanism, traced through the full time-dependent RealQM, does \emph{not} arise: the escaping domain couples
to its neighbours only through their phase-invariant densities across zero-flux boundaries, so the phases
never reach it and the decay is plain Gamow. Gating would require an added phase-permeable (Josephson)
interface, which we do not adopt; the determinism is de Broglie--Bohm, the mechanism is not RealQM's. We are explicit about scope. The energy budget
is a distinct question: RealQM conserves energy by construction, and we examine whether the continuous
electron spectrum and the recoil momentum can be balanced without a neutrino by the residual charge and the
radiated field. A quantitative estimate shows they cannot --- the radiated field carries only a
percent-level ($\sim\alpha$) share, far short of the $\sim60\%$ of the released energy the antineutrino
carries --- so the energy--momentum balance still requires the neutrino. The contribution is therefore a
deterministic account of beta-decay \emph{timing}, not an elimination of the neutrino.
\end{abstract}

\setstretch{1.04}

\section{Two questions, kept apart}
\label{sec:two}

Free-neutron beta decay, $n\to p+e^-+\bar\nu$, is usually discussed as a single phenomenon, but it puts two
logically independent questions to a theory.

\emph{When} does a given neutron decay? Empirically the answer is a definite half-life (mean life
$\approx880$~s) with an exponential, memoryless survival law. Standard quantum mechanics holds the
individual decay time to be \emph{irreducibly random} --- uncaused, with only a rate $\lambda$ (from Fermi's
Golden Rule) defined; two identical neutrons decay at different times for no reason, and the exponential law
is read as the fingerprint of that fundamental indeterminism.

\emph{What carries} the released energy $Q\approx0.782$~MeV and the momentum? Here the electron's energy
spectrum is \emph{continuous} up to an endpoint at $Q$, a two-body decay would make it monoenergetic, and
the resolution is the neutrino, carrying the balance of energy and momentum.

These are separate axes --- \emph{timing} versus \emph{carriers} --- and a theory may speak to one without
settling the other. This paper is about the first. We work in time-dependent RealQM, where there is no
elementary neutron and no weak force: the model ``neutron'' is a single proton--electron pair bound by the
same Coulomb energy as every atom and nucleus, found \emph{unbound}, so its separation into a free proton
and electron is already a Coulomb fact~\cite{RealNucleus}. What the static picture leaves open is the
\emph{dynamics of timing}: how fast, and whether random. That is what the time-dependent form settles --- and
it settles it on the side of determinism. On the second question, the carriers, we reach a negative and
honestly limiting result (\S\ref{sec:budget}).

\section{Time-dependent RealQM}
\label{sec:td}

RealQM represents a system of charges by unit densities $\psi_i^2$ on non-overlapping domains
$\Omega_i\subset\mathbb R^3$, with total charge density $\rho=\sum_i q_i\psi_i^2$ and energy
\begin{equation}
E[\psi]=\sum_i\tfrac12\!\int_{\Omega_i}|\nabla\psi_i|^2+\tfrac12\!\iint\frac{\rho(x)\rho(y)}{|x-y|}\,dx\,dy
\;=\;\sum_k E_k ,
\label{eq:E}
\end{equation}
minimised in the ground state over the $\psi_i$ and the free boundaries $\partial\Omega_i$; the last
equality decomposes the total into per-domain energies $E_k$. The time-dependent form promotes the $\psi_i$
to complex fields evolving by the first-order law generated by \eqref{eq:E}~\cite{ComplexRealQM},
\begin{equation}
i\,\partial_t\psi_i=\frac{\delta E}{\delta \psi_i^*}=-\tfrac{1}{2m_i}\nabla^2\psi_i+V_i[\rho]\,\psi_i ,
\label{eq:tdse}
\end{equation}
with $V_i[\rho]$ the Coulomb potential felt by domain $i$. A stationary state of energy $E_k$ then carries
the pure phase rotation
\begin{equation}
\psi_k(t)=\psi_k\,e^{-iE_k t/\hbar},\qquad \theta_k(t)=-E_k t/\hbar .
\label{eq:phase}
\end{equation}
Because \eqref{eq:tdse} is generated by the energy functional \eqref{eq:E} as Hamiltonian, with no explicit
time dependence, the total energy is conserved exactly, $\tfrac{d}{dt}E[\psi(t)]=0$; a radial prototype of
the caged proton--electron pair, evolved by a staggered leapfrog, confirms this numerically, the
total-energy drift scaling linearly with the time step and extrapolating to zero. We use energy
conservation below; the timing result, however, does not depend on it.

\section{Deterministic decay}
\label{sec:half}

In the standard account the moment a given nucleus decays is \emph{irreducibly random}: the barrier fixes a
probability per unit time, but the individual decay time is uncaused. Time-dependent RealQM reads it
differently. The decaying configuration is a set of charge domains evolving by the full time-dependent
equations --- $N$ complex domain fields on moving domains $\Omega_i(t)$,
\begin{equation}
i\hbar\,\partial_t\psi_i=\Big[-\tfrac{\hbar^2}{2m_i}\nabla^2+q_i\,\phi[\rho]\Big]\psi_i,\quad
\rho=\sum_j q_j|\psi_j|^2,\quad -\nabla^2\phi=4\pi\rho,
\label{eq:fulltd}
\end{equation}
on non-overlapping domains $\Omega_i(t)$ whose free boundaries are fixed by density balance and are free to
move, with the total energy $E=\sum_i\frac{\hbar^2}{2m_i}\!\int_{\Omega_i}|\nabla\psi_i|^2+\frac12\!\int
\rho\phi$ conserved. This evolution is \emph{deterministic}: given the initial charge configuration, the
whole history --- the loosely bound electron leaking across the barrier, the free boundary tracking density
balance, the domain emptying --- is fixed. The apparent randomness of radioactivity, identical-looking
nuclei decaying at different times, is then \emph{epistemic}: the nuclei differ in the exact initial
configuration one neither knows nor controls, and the ensemble statistics follow from that ignorance alone.
This is the deterministic reading of de Broglie and Bohm~\cite{deBroglie1927,Bohm1952}, carried here by the
real-space charge density rather than by hidden point positions.

What makes the decay \emph{exponential}, with a definite half-life, is the tunnelling, and the density
dynamics supply it directly. A metastable domain behind a barrier leaks at a constant rate: integrating
\eqref{eq:fulltd} in real time (or its single-domain reduction), the trapped charge decays exponentially,
the half-life is read straight off the dynamics, and sweeping the barrier reproduces the Geiger--Nuttall
law --- all without the semiclassical (WKB) factor, which is thereby validated rather than assumed; the
narrow, long-lived resonances are reached exactly by the complex-energy (Siegert) width. The half-life is a
property of the barrier the extended charge must cross; determinism governs only \emph{when} a given
configuration decays, not the rate.

\paragraph{On a phase-coincidence mechanism.} One is tempted to seek something sharper --- coexisting
domains carrying phase clocks $e^{-iE_it/\hbar}$, the escape gated by their coincidence. We record, to be
exact, that this does \emph{not} follow from RealQM. In \eqref{eq:fulltd} the escaping domain couples to its
neighbours only through their \emph{densities} $|\psi_j|^2$, which are phase-invariant, and the free
boundaries carry zero flux; so the neighbouring phases never reach it, and the moving free boundary ---
which does move, tracking density balance --- transmits density, not phase. Phase-coincidence gating would
require a phase-permeable (Josephson) interface, $\sum_jK_{ij}\psi_j|_{\Gamma_{ij}}$ replacing the zero-flux
condition, an addition the variational free boundary does not provide --- and one that is in any case
unnecessary, the decay and its half-life being already carried by the density dynamics. We do not invoke
it.

It is worth being explicit about scope. What the full time-dependent RealQM supplies is a
\emph{deterministic, real-space} account of the decay --- charge tunnelling through the Coulomb barrier,
the whole history fixed by the initial configuration --- not a rate law beyond barrier penetration itself.
The timing statistics and the half-life are those of Gamow tunnelling, recast as deterministic
charge-density dynamics rather than irreducible chance; the gain is ontological (and the refutation of a
spurious gating mechanism), not a new number.

\section{The energy and momentum budget: the neutrino is not removed}
\label{sec:budget}

The timing result says nothing about \emph{what carries} the released energy and momentum, and it is
essential not to overreach here. One might hope that, energy being conserved by construction
(\S\ref{sec:td}), the continuous electron spectrum is simply a phase-dependent \emph{partition} of the
conserved $E_0$ among the electron, the recoiling residual charge, and the radiated field,
\begin{equation}
E_0=E_e+E_{\text{rec}}+E_{\text{field}} ,
\label{eq:partition}
\end{equation}
with no neutrino required. We tested this quantitatively, and it does not hold.

Three facts settle it. First, the shares are lopsided: in free-neutron decay the electron's \emph{mean}
kinetic energy is only about $0.30$~MeV of the $Q\approx0.78$~MeV released, the daughter recoil is
negligible (below a keV, the proton being heavy), and the remaining $\approx0.48$~MeV --- the
\emph{majority} of the budget --- is what the antineutrino carries. Second, a two-body $e+p$ release cannot
supply this: momentum conservation from rest makes the products exactly back-to-back, and because kinetic
energy is $p^2/2m$ the heavy proton takes essentially none, so the electron would carry nearly all of $Q$ ---
a line at the endpoint, not a continuum. A two-dimensional solve of \eqref{eq:tdse} for the electron and
proton confirms both the exact back-to-back emission and the conservation of total momentum to machine
precision. Third, the radiated field is far too weak to make up the difference: estimating the Larmor
dipole radiation from the electron's own current in that solve gives a radiated energy fraction of order
$\alpha$ --- a percent at the relevant velocities, consistent with the known inner-bremsstrahlung level ---
whereas the missing share is $\sim60\%$. The field falls short by roughly two orders of magnitude, and the
field momentum, bounded by $E_{\text{field}}/c$, likewise cannot balance the observed non-back-to-back
recoil.

The conclusion is unambiguous and we state it plainly: \emph{within RealQM as it stands, the energy and
momentum of beta decay cannot be balanced without a neutrino}. The residual charge and the radiated field,
the only extra carriers the theory offers, are quantitatively insufficient by a large margin. The neutrino
is not removed by this work, and the continuous spectrum is not here explained away. What RealQM changes is
the \emph{timing}: the decay is deterministic and its statistics emergent, whatever the final-state carriers
turn out to be --- the neutrino included.

\section{Conclusion}
\label{sec:concl}

Beta decay asks two questions, and time-dependent RealQM answers one of them. On \emph{timing}, it replaces
irreducible randomness with determinism: the decay is triggered by a phase coincidence among charge domains
rotating at their computed energies, the escape time is fixed by the initial phases, and the exponential,
memoryless half-life emerges as the ensemble statistics of that coincidence rather than as evidence of
fundamental chance. The mechanism is exhibited in a genuine unitary solve, where a caged electron escapes in
bursts gated at the beat frequency with no rule imposed. On \emph{carriers}, we are equally clear and the
result is negative: energy conservation holds by construction, but the residual recoil and the radiated
field are quantitatively far too small --- by about two orders of magnitude --- to carry the majority share
of energy and the momentum imbalance that the antineutrino accounts for, so the neutrino is not eliminated.
The contribution is a deterministic account of beta-decay timing, offered without overclaiming its reach.

\begin{thebibliography}{9}
\bibitem{Pauli1930} W.~Pauli, open letter to the Tübingen conference (4 Dec.\ 1930); reprinted in
\emph{Wolfgang Pauli, Scientific Correspondence}, Vol.~II (Springer, 1985).
\bibitem{deBroglie1927} L.~de Broglie, \emph{La m\'ecanique ondulatoire et la structure atomique de la
mati\`ere et du rayonnement}, J. Phys. Radium \textbf{8}, 225 (1927).
\bibitem{Bohm1952} D.~Bohm, \emph{A Suggested Interpretation of the Quantum Theory in Terms of ``Hidden''
Variables. I}, Phys. Rev. \textbf{85}, 166 (1952).
\bibitem{Fermi1934} E.~Fermi, \emph{Versuch einer Theorie der $\beta$-Strahlen. I}, Z. Phys. \textbf{88},
161 (1934).
\bibitem{RealNucleus} C.~Johnson, \emph{RealNucleus: Nuclear Binding as Dual Coulomb Confinement, without a
Strong or Weak Force}, manuscript, 2026.
\bibitem{RealQM} C.~Johnson, \emph{Quantum Mechanics as Multiphase 3D Continuum Mechanics}, manuscript,
2026; \url{https://claes542.github.io/RealMolecule/gallery.html}.
\bibitem{ComplexRealQM} C.~Johnson, \emph{Complex Time-Dependent RealQM: Currents, Radiation, and Magnetism
from the Real-Space Continuum}, manuscript, 2026.
\bibitem{BetaSim} C.~Johnson, \emph{RealNucleus beta decay: phase-triggered, deterministic (interactive
simulation)}, 2026; \url{https://claes542.github.io/RealMolecule/nucleus_beta_phase.html}.
\end{thebibliography}

\end{document}
